\chapter{Tetanus}
\index{Tetanus}

\medskip
\noindent\textit{Educational reference; always seek professional advice for individual care.}

\section*{Definition and Overview}
Tetanus is a serious, potentially fatal bacterial infection of the nervous system, caused by a toxin produced by the bacterium Clostridium tetani. The organism is found in soil, dust, and animal and human faeces, and it enters the body through a wound, usually a puncture, cut, burn, or bite contaminated with soil or dirt. Once inside, the bacteria multiply in damaged tissue and release a toxin that travels to the spinal cord and brain, where it blocks the signals that normally relax muscles. The result is uncontrollable, painful muscle spasms and stiffness, beginning classically in the jaw, which is why tetanus is often called lockjaw, and then spreading to the neck, chest, back, and limbs. In severe cases the spasms interfere with breathing and with the functioning of the heart and other organs, and the disease can be fatal even with the best hospital care. The critical, hopeful fact is that tetanus is entirely preventable: a safe and effective vaccine exists, it is part of the routine childhood immunisation schedule in many countries, and boosters are recommended for adults throughout life. Cases in many countries are now very rare, and almost all occur in people who are unvaccinated, under-vaccinated, or overdue for a booster, most commonly older adults. This chapter explains how tetanus is caught, what it looks like, how it is treated, and, most importantly, how vaccination protects against it.

\section*{Epidemiology}
Tetanus is extremely rare in many countries and other countries with high vaccination coverage, but it remains a major cause of death globally, particularly in low-income regions where immunisation is incomplete. In many countries, only a handful of cases are reported each year, and deaths are rarer still, whereas worldwide the World Health Organization estimates that hundreds of thousands of cases occur annually, with the highest burden in South Asia and Africa and in the newborn. The pattern of who is affected is instructive: the disease now occurs almost exclusively in people who are not fully vaccinated, including older adults whose childhood protection has waned and who have not had a booster, migrants who missed childhood vaccination, people who inject drugs into skin or muscle, and occasional travellers returning from countries with poor coverage. Occupational and leisure exposure matters too, because farmers, gardeners, and outdoor workers are more likely to get the soil-contaminated wounds that introduce the organism. Neonatal tetanus, once common, has been virtually eliminated in many countries through maternal vaccination and clean delivery practices. Every local case represents a missed opportunity for prevention.

\section*{Aetiology and Pathophysiology}
Clostridium tetani is a hardy bacterium that exists as spores, able to survive for years in soil and dust. The spores do not cause illness on their own; they germinate only in wounds where oxygen is low, such as deep puncture wounds, wounds with dead tissue, burns, and compound fractures.

\begin{itemize}
    \item \textbf{Spore entry through a wound:} the bacterium enters through any break in the skin, from a nail, splinter, cut, burn, bite, or scratch, particularly wounds contaminated with soil, dust, or faeces
    \item \textbf{Germination in damaged tissue:} spores germinate in wounds with poor blood supply, dead tissue, or foreign material, where oxygen levels are low
    \item \textbf{Toxin production:} the growing bacteria release tetanospasmin, the toxin responsible for the disease
    \item \textbf{Spread to the nervous system:} the toxin travels along nerve pathways to the spinal cord and brainstem
    \item \textbf{Blocking of inhibitory signals:} the toxin stops the normal chemical signals that relax muscles after contraction, leaving excitatory signals unopposed
    \item \textbf{Sustained muscle spasm:} with relaxation blocked, muscles contract continuously, producing the rigidity, spasms, and lockjaw of tetanus
    \item \textbf{Autonomic disruption:} in severe disease, the toxin also disturbs the automatic nervous system, causing instability of heart rate and blood pressure
\end{itemize}

The incubation period is usually three to twenty-one days after the wound, and the shorter the incubation period, the more severe the disease tends to be. The toxin acts so strongly that the disease can progress even after the wound itself has healed.

\section*{Clinical Features}
\begin{itemize}
    \item \textbf{Trismus (lockjaw):} stiffness and spasm of the jaw muscles, so that the mouth becomes difficult or impossible to open; this is the classic first symptom
    \item \textbf{Risus sardonicus:} a fixed, grimacing smile caused by spasm of the facial muscles
    \item \textbf{Neck and back stiffness:} rigidity spreads to the neck, and in severe cases to the back, sometimes arching the body backwards (opisthotonos)
    \item \textbf{Painful generalised spasms:} sudden, intense, whole-body muscle spasms, sometimes triggered by noise, touch, light, or any startle
    \item \textbf{Difficulty swallowing:} spasm of the throat muscles makes swallowing painful and difficult
    \item \textbf{Chest and abdominal rigidity:} stiffness of the chest and abdominal muscles can impair breathing and is a danger sign
    \item \textbf{Fever and sweating:} a moderate fever and sweating are common
    \item \textbf{Autonomic instability:} in severe cases, fluctuating blood pressure and fast or irregular heart rhythm
    \item \textbf{Consciousness is usually preserved:} unlike many neurological emergencies, the person remains alert, which makes the experience especially distressing
\end{itemize}

The disease is graded from mild, with localised or mild generalised rigidity, to severe, with frequent generalised spasms and breathing or autonomic complications.

\section*{Differential Diagnosis}
\begin{itemize}
    \item \textbf{Trismus from dental or jaw problems:} dental abscess, temporomandibular joint dysfunction, or tonsillitis can limit jaw opening, but they are accompanied by local pain and swelling, not generalised spasm
    \item \textbf{Drug-induced muscle rigidity:} antipsychotic medicines can cause acute muscle stiffness and spasm, but the pattern and history differ
    \item \textbf{Strychnine poisoning:} causes similar generalised spasms, distinguished by history of exposure and a rapid onset
    \item \textbf{Hypocalcaemia:} low blood calcium can cause muscle twitching, cramping, and spasm, usually with other metabolic signs
    \item \textbf{Brainstem or neurological disease:} some strokes, encephalitis, or multiple sclerosis can cause jaw or neck stiffness, but with other neurological findings
    \item \textbf{Epilepsy or psychogenic spasms:} seizures or non-epileptic attacks cause different patterns and usually alter consciousness
    \item \textbf{Rabies:} an extremely rare viral infection with swallowing spasms, agitation, and hydrophobia, considered if there has been animal exposure
\end{itemize}

Because tetanus is so rare, it can be missed early when the only symptom is jaw stiffness; a history of a recent wound and incomplete vaccination should immediately raise the suspicion.

\section*{When to Seek Medical Care}
Tetanus is a medical emergency that requires hospital treatment. Anyone with suspected tetanus must be assessed urgently, and any wound in a person whose vaccination is not up to date should be assessed promptly for a booster or, if indicated, tetanus immune globulin.

Seek medical care promptly if you have:
\begin{itemize}
    \item A wound contaminated with soil, dirt, rust, or animal or human faeces, and you are unsure whether your tetanus vaccination is up to date
    \item A deep puncture wound, a wound with dead tissue, a burn, or an animal bite, particularly if your last booster was more than ten years ago
    \item Stiffness or spasm of the jaw, difficulty opening the mouth, or unexplained muscle stiffness, especially after a recent wound
\end{itemize}

\textbf{Contact your local emergency services for:}
\begin{itemize}
    \item Stiffness or spasm spreading from the jaw to the neck, back, or chest, or difficulty breathing or swallowing
    \item Sudden, painful, whole-body spasms
    \item Fever with rigidity of the neck or back, or any suspicion of tetanus in an unvaccinated or under-vaccinated person
\end{itemize}

\section*{Diagnosis and Investigations}
The diagnosis of tetanus is clinical. There is no blood test that confirms it in routine practice, and the diagnosis rests on the history of a wound and the characteristic muscle rigidity and spasms.

\begin{itemize}
    \item \textbf{Clinical assessment:} the combination of trismus, generalised rigidity, and painful spasms, with a history of a recent wound and incomplete vaccination, is the basis of the diagnosis
    \item \textbf{Wound history:} establishing the nature, timing, and contamination of the wound, and the vaccination status of the person
    \item \textbf{Vaccination history:} checking childhood records and booster dates is central to both diagnosis and prevention
    \item \textbf{Exclusion of mimics:} blood tests for calcium, and review of medicines, exclude metabolic or drug causes of spasm
    \item \textbf{Culture:} occasionally the wound is swabbed for the bacterium, but culture is often negative and does not rule tetanus in or out
    \item \textbf{Blood tests:} used to assess complications such as dehydration, infection, or damage to organs in severe disease
    \item \textbf{Monitoring:} severe cases are managed in intensive care with continuous monitoring of breathing, heart rate, and blood pressure
\end{itemize}

\section*{Management}
Tetanus is treated in hospital, and severe cases are managed in an intensive care unit. The aims are to remove the source of toxin, neutralise toxin that has not yet entered the nerves, control the spasms, and support the body while the effect of the toxin wears off.

\begin{itemize}
    \item \textbf{Wound care:} the wound is thoroughly cleaned and dead tissue is removed, removing the environment where the bacteria grow
    \item \textbf{Tetanus immune globulin:} given as soon as possible to neutralise free toxin in the bloodstream; it does not reverse toxin already bound to nerves
    \item \textbf{Antibiotics:} metronidazole (or penicillin) is given to kill the bacteria in the wound
    \item \textbf{Sedation and muscle relaxants:} medicines such as benzodiazepines and other relaxants control the spasms and rigidity
    \item \textbf{Intensive care support:} severe cases may require breathing support with a ventilator, and careful control of blood pressure and heart rhythm
    \item \textbf{Nursing care:} quiet surroundings, gentle handling, and measures to prevent pressure sores and complications of immobility
    \item \textbf{Nutrition:} feeding support, often through a tube, maintains strength during a prolonged illness
    \item \textbf{Pain and comfort:} analgesia and sedation are important, because the spasms are intensely painful and frightening
    \item \textbf{Vaccination:} every person who has tetanus is given the full vaccine course after recovery, because the disease itself does not confer immunity
\end{itemize}

Recovery is slow, typically taking weeks to months, because the toxin must be cleared from the nervous system and damaged muscle tissue repaired.

\section*{Complications}
\begin{itemize}
    \item \textbf{Respiratory failure:} spasm of the chest and airway muscles, or spasm of the voice box, can stop breathing and is the commonest cause of death
    \item \textbf{Cardiac and autonomic instability:} severe tetanus causes dangerous swings in blood pressure and heart rhythm, which can be fatal
    \item \textbf{Fractures and muscle injury:} violent spasms can fracture vertebrae and tear muscles
    \item \textbf{Secondary infections:} pneumonia, sepsis, and wound infection complicate the course in hospital
    \item \textbf{Complications of prolonged ICU care:} deep vein thrombosis, pressure sores, and the effects of long-term ventilation
    \item \textbf{Long recovery:} survivors can take weeks to months to recover, with muscle weakness and stiffness persisting
\end{itemize}

\section*{Prognosis and Outcomes}
The outcome of tetanus depends mainly on the severity of the disease and the quality of supportive care. In many countries and other well-resourced countries, with intensive care treatment, the outlook has improved dramatically, but tetanus remains a dangerous disease: even in good hospitals, a meaningful proportion of severe cases are fatal, and in countries without intensive care the death rate is far higher. The disease is worse in the very young, the very old, and those with rapid onset of symptoms. For survivors, recovery is slow and often complete, although some people are left with stiffness, weakness, or the psychological aftermath of a frightening illness. The most important fact, however, is preventive: tetanus is almost entirely avoidable. A full course of vaccination followed by boosters provides protection for decades, which is why the great majority of people in many countries will never know the disease. No death from tetanus should ever occur in a country with safe and effective vaccines, and the tragedy of the rare local case is that it was preventable.

\section*{Prevention and Self-Care}
\begin{itemize}
    \item Complete the routine childhood vaccination schedule, which includes tetanus, and check with a GP or on your immunisation record that you are up to date
    \item Adults should have a tetanus booster at the recommended intervals, and a booster is also given when a tetanus-prone wound occurs and it has been more than ten years since the last dose
    \item Clean any wound thoroughly with clean water and antiseptic, and remove any visible dirt or foreign material
    \item Seek prompt assessment for deep puncture wounds, wounds contaminated with soil or faeces, burns, and animal bites, especially if vaccination status is unclear
    \item Wear protective gloves and footwear when gardening, farming, or handling soil
    \item Pregnant women should have the recommended boosters, which protect both mother and newborn
    \item Travellers to countries with limited healthcare should confirm their tetanus vaccination is current before departure
    \item People who inject drugs should be up to date with vaccination and seek help for any skin or soft-tissue infection
\end{itemize}

\section*{Sources and Further Reading}
The following organisations provide reliable information about tetanus and its prevention.

\begin{itemize}
    \item Australian Department of Health and Aged Care. \textit{Tetanus and the national immunisation program.} https://www.health.gov.au/
    \item healthdirect Australia. \textit{Tetanus.} https://www.healthdirect.gov.au/
    \item Australian Immunisation Handbook. \textit{Tetanus vaccine information.} https://immunisationhandbook.health.gov.au/
    \item World Health Organization. \textit{Tetanus fact sheet.} https://www.who.int/
    \item NHS. \textit{Tetanus.} https://www.nhs.uk/
    \item Mayo Clinic. \textit{Tetanus: symptoms, causes, and prevention.} https://www.mayoclinic.org/
\end{itemize}

\clearpage
